\section{Kết quả}\label{sec:ket-qua}

\subsection{Nhận diện: tầng xác minh cắt hai phần ba báo nhầm}

Bảng~\ref{tab:tang} và Hình~\ref{fig:tang} là kết quả chấm một lần trên 198 khung test, qua đúng đường IOT-CV chạy: phát lại bản ghi, ghi CSV từng tầng, chấm bằng \texttt{xla.eval.detection}. So YOLO11s một tầng với YOLO11s + verifier EfficientNet-B0 + lọc track (ba tầng) trên cùng 198 khung: mAP50 tăng từ 0,120 lên 0,210, hiệu $+0{,}090$ \ci{$+0{,}026$}{$+0{,}189$}, và báo nhầm trên khung bay không người giảm từ 0,881 xuống 0,289 hộp/khung, hiệu $-0{,}591$ \ci{$-1{,}033$}{$-0{,}273$}. Cả hai CI không chạm 0. Recall ở ngưỡng vận hành giữ nguyên 0,396 (59/149 người): verifier không bỏ người nào mà detector đã bắt được. Trên 300 khung trong nhà, báo nhầm ở mức 0,003 hộp/khung.

\begin{table}[!htbp]
\centering\footnotesize
\caption[Ba tầng nhận diện trên 198 khung test thật]{Ba tầng nhận diện trên 198 khung test thật (39 khung có người, 149 người). CI 95\,\% bootstrap theo khối 30 s. Báo nhầm tính trên 159 khung bay không người và 300 khung trong nhà.}
\label{tab:tang}
\setlength{\tabcolsep}{3.5pt}
\begin{tabular}{@{}l c c c c c@{}}
\toprule
Tầng & mAP50 [CI 95\,\%] & Recall & Báo nhầm, bay [CI 95\,\%] & Trong nhà & Ngưỡng \\
\midrule
1 detector & 0,119 \ci{0,032}{0,197} & 0,396 & 0,874 \ci{0,333}{1,486} & 0,003 & 0,077 \\
2 + verifier & 0,210 \ci{0,122}{0,312} & 0,396 & 0,289 \ci{0,036}{0,631} & 0,003 & 0,350 \\
3 + track $k1/n1$ & 0,210 \ci{0,122}{0,312} & 0,396 & 0,289 \ci{0,036}{0,631} & 0,003 & 0,350 \\
\midrule
ba tầng $-$ một tầng & $+0{,}090$ \ci{$+0{,}026$}{$+0{,}189$} & $+0{,}000$ & $-0{,}591$ \ci{$-1{,}033$}{$-0{,}273$} & & \\
\bottomrule
\end{tabular}
\end{table}

\begin{figure}[!htbp]
\centering
% Ba tầng trên 198 khung test (bảng t7, system/score.json). Thanh lỗi: CI 95 % bootstrap theo khối 30 s.
\begin{subfigure}[b]{0.49\textwidth}\centering
\begin{tikzpicture}
\begin{axis}[baocao, width=7.2cm, height=5.0cm, ybar, bar width=16pt,
  ymin=0, ymax=0.40, ylabel={mAP50 test}, xmin=0.4, xmax=3.6, xtick={1,2,3},
  xticklabels={{tầng 1\\detector},{tầng 2\\+ verifier},{tầng 3\\+ track}},
  x tick label style={align=center,font=\scriptsize}, /pgf/number format/precision=2,
  error bars/y dir=both, error bars/y explicit, error bars/error bar style={line width=0.7pt,black!70},
  error bars/error mark options={rotate=90,mark size=3pt,black!70}]
\addplot[fill=cCam, draw=none, bar shift=0pt] coordinates {(1,0.119) += (0,0.078) -= (0,0.087)};
\addplot[fill=cTim!85, draw=none, bar shift=0pt] coordinates {(2,0.210) += (0,0.102) -= (0,0.088)};
\addplot[fill=cLuc!85, draw=none, bar shift=0pt] coordinates {(3,0.210) += (0,0.102) -= (0,0.088)};
\node[font=\scriptsize, anchor=south] at (axis cs:1,0.205) {0,119};
\node[font=\scriptsize, anchor=south] at (axis cs:2,0.320) {0,210};
\node[font=\scriptsize, anchor=south] at (axis cs:3,0.320) {0,210};
\end{axis}
\end{tikzpicture}
\caption{mAP50 trên 198 khung test}
\label{fig:tang-map}
\end{subfigure}\hfill
\begin{subfigure}[b]{0.49\textwidth}\centering
\begin{tikzpicture}
\begin{axis}[baocao, width=7.2cm, height=5.0cm, ybar, bar width=16pt,
  ymin=0, ymax=1.75, ylabel={hộp sai / khung bay trống}, xmin=0.4, xmax=3.6, xtick={1,2,3},
  xticklabels={{tầng 1\\detector},{tầng 2\\+ verifier},{tầng 3\\+ track}},
  x tick label style={align=center,font=\scriptsize}, /pgf/number format/precision=2,
  error bars/y dir=both, error bars/y explicit, error bars/error bar style={line width=0.7pt,black!70},
  error bars/error mark options={rotate=90,mark size=3pt,black!70}]
\addplot[fill=cCam, draw=none, bar shift=0pt] coordinates {(1,0.874) += (0,0.612) -= (0,0.541)};
\addplot[fill=cTim!85, draw=none, bar shift=0pt] coordinates {(2,0.289) += (0,0.342) -= (0,0.253)};
\addplot[fill=cLuc!85, draw=none, bar shift=0pt] coordinates {(3,0.289) += (0,0.342) -= (0,0.253)};
\node[font=\scriptsize, anchor=south] at (axis cs:1,1.50) {0,874};
\node[font=\scriptsize, anchor=south] at (axis cs:2,0.64) {0,289};
\node[font=\scriptsize, anchor=south] at (axis cs:3,0.64) {0,289};
\end{axis}
\end{tikzpicture}
\caption{Báo nhầm trên 159 khung bay không người}
\label{fig:tang-fp}
\end{subfigure}

\caption[Ba tầng nhận diện trên 198 khung test]{Ba tầng nhận diện trên 198 khung test, thanh lỗi là CI 95\,\% bootstrap theo khối. Phần tăng mAP50 và phần giảm báo nhầm đều đến từ tầng 2.}
\label{fig:tang}
\end{figure}

Tầng 1 chạy qua IOT-CV cho đúng kết quả của detector chấm bằng XLA (hiệu mAP50 bằng 0, lệch một hộp báo nhầm trên 159 khung), nên phần chênh ở tầng 2 là của verifier. Trên \texttt{cal}, verifier ở $\tau=0{,}35$ bỏ 13/15 hộp sai; trên \texttt{devpool} (dự đoán ngoài fold) nó bỏ 25/79 hộp sai mà vẫn giữ đủ 58/218 người. Tầng 3 ở $k=1$ trùng tầng 2 từng hộp, đúng như luật D7 chọn: lọc track $k\geq2$ bỏ nốt 2 hộp sai trên \texttt{cal} nhưng mất ít nhất 10 người trên \texttt{devpool}. Ba gói ONNX (nhận diện, xác minh, làm rõ) đều qua kiểm golden của IOT-CV trên cả CUDA và CPU (detector IoU $\geq$ 0,9998, verifier lệch điểm $\leq$ 0,0003, bộ lọc cổ điển lệch $\leq$ 2 mức).

\subsection{Làm rõ ảnh: bộ lọc cổ điển thắng ảnh gốc 25/25}

Trên 198 khung test, bộ lọc cổ điển ONNX hạ tỉ lệ điểm hồng từ 0,943\,\% xuống 0,102\,\% (giảm 89,2\,\%) và hạ dark channel trung vị từ 139,0 xuống 56,5 (giảm 59,4\,\%); trên 50 khung dev các số tương ứng là 91,4\,\% và 60,4\,\%, nên dev và test không lệch nhau. Hình~\ref{fig:enh} là hai khung bay thật trước và sau: lớp sương xám được gỡ, cỏ về màu xanh, vệt hồng trên đường đất nhạt đi, người và bóng người nổi rõ khỏi nền.

\begin{figure}[!htbp]
\centering
\foreach \i/\c in {0/{\texttt{002\_00428}},2/{\texttt{007\_00765}}}{%
\begin{subfigure}[b]{0.285\textwidth}\centering\includegraphics[width=\linewidth]{fig/anh\i_raw.jpg}\caption*{gốc, \c}\end{subfigure}\hfill
\begin{subfigure}[b]{0.285\textwidth}\centering\includegraphics[width=\linewidth]{fig/anh\i_classic.jpg}\caption*{cổ điển ONNX}\end{subfigure}\hfill
\begin{subfigure}[b]{0.205\textwidth}\centering\includegraphics[width=\linewidth]{fig/anh\i_raw_zoom.jpg}\caption*{gốc, 2$\times$}\end{subfigure}\hfill
\begin{subfigure}[b]{0.205\textwidth}\centering\includegraphics[width=\linewidth]{fig/anh\i_classic_zoom.jpg}\caption*{cổ điển, 2$\times$}\end{subfigure}\par\smallskip}
\caption[Khung bay thật qua bộ lọc cổ điển ONNX]{Khung bay thật ngày 02/11 qua bộ lọc cổ điển ONNX: ảnh ra là đầu ra ORT đã lưu làm golden của gói ONNX làm rõ. Hai cột phải phóng 2$\times$ vùng có người (góc trên trái).}
\label{fig:enh}
\end{figure}

Bảng~\ref{tab:sosanh} đặt bộ lọc cổ điển cạnh các phương án đã đo. Ở thử mù A/B lượt cả khung (Hình~\ref{fig:ab}), bộ lọc cổ điển thắng ảnh gốc 25/25, Wilson \ci{0,867}{1,000}; năm cặp lặp đổi bên đều được chọn như lần đầu. Real-ESRGAN x2 tinh chỉnh cũng thắng ảnh gốc 25/25, còn khi đặt cạnh bộ lọc cổ điển nó được chọn 16/25, Wilson \ci{0,445}{0,798}: người chấm không phân biệt được hai cách ở cỡ cả khung. Hàng rào báo nhầm của ảnh offline qua (detector trên 150 khung \texttt{cal}: 15 hộp ở ảnh gốc, 10 hộp ở ảnh offline, hiệu $-0{,}033$ \ci{$-0{,}123$}{$+0{,}039$}).

\begin{table}[!htbp]
\centering\footnotesize
\caption[So sánh các phương pháp làm rõ ảnh]{So sánh các phương pháp làm rõ ảnh. Thời gian đo trên RTX 4060 Laptop, ONNX Runtime 1.24.4 trừ khi ghi khác; ``chuyến'' là một chuyến bay 20 phút (6\,750 khung ở 5,62 khung/s).}
\label{tab:sosanh}
\begin{tabularx}{\textwidth}{@{}>{\raggedright\arraybackslash}p{2.7cm} >{\raggedright\arraybackslash}p{3.3cm} >{\raggedright\arraybackslash}p{3.0cm} X@{}}
\toprule
Phương pháp & Kết quả đo & Tốc độ & Lý do chọn hay không chọn \\
\midrule
Bộ lọc cổ điển ONNX & A/B thắng gốc 25/25; hồng $-$89,2\,\%, dark channel $-$59,4\,\% (test) & 19,3 / 25,3 ms p50/p95 CUDA; 29,3 / 30,6 ms CPU & chọn, bật sẵn: không học nên không mang thiên lệch của ảnh huấn luyện, khử mù theo chính khung đang xem \\
\addlinespace
NAFNet-w16 & dark channel $-$36,8\,\%, màu ra nguội & 20,61 / 27,89 ms (cùng lượt đo) & không chọn: khử mù yếu hơn, chậm hơn 2,55 ms p95 \\
\addlinespace
Lai cổ điển $\to$ NAFNet & hết hồng; dark channel trung vị 84,0 đến 111,5 (cổ điển 56,0), dev & 7\,600 bước, 25 phút GPU & không xuất: mạng học theo đích VisDrone sạch nên trả lại màn mờ bộ lọc đã gỡ \\
\addlinespace
Real-ESRGAN x2 tinh chỉnh (offline) & A/B thắng gốc 25/25; với cổ điển 16/25 & 0,404 s/khung, 0,758 giờ GPU/chuyến & không chọn: A/B cả khung không phân biệt được với cổ điển, trong khi chậm hơn khoảng 20 lần \\
\addlinespace
SwinIR-M x2 & & 2,70 s/khung PyTorch; 17,4 s ORT & dừng ở bench: một chuyến cần 32,6 giờ GPU \\
\addlinespace
BasicVSR++ (5 khung) & & 3,34 s/khung PyTorch & dừng ở probe: ORT 1.24.4 không có kernel \texttt{DeformConv} \\
\bottomrule
\end{tabularx}
\end{table}

\subsection{Chạy trong phần mềm drone}

IOT-CV nạp ba gói ONNX và phát lại 600 khung của chuyến 02/11 (khung 284 đến 884) đúng nhịp thật, so hai cấu hình trong một tiến trình, bốn lượt xen kẽ ABBA: nhận diện một tầng + NAFNet (A) và ba tầng + bộ lọc cổ điển (B). Ở luồng mặc định có người xem ``Ảnh rõ'', p95 tổng của B là 89,5 ms so với 99,1 ms của A (p50 76,3 so với 84,5 ms), và tỉ lệ khung bị bỏ là 4,2 và 4,7\,\% so với 7,8 và 8,2\,\% (Hình~\ref{fig:ms}). Bộ lọc cổ điển nhẹ hơn NAFNet khoảng 15 ms p95 (19,3 so với 34,4 ms), đủ bù phần verifier thêm vào (p95 khoảng 9 ms, p50 0,03 ms vì phần lớn khung không có hộp qua tầng 1). Lệch giữa hai lượt cùng cấu hình $\leq$ 3,4 ms, nhỏ hơn chênh giữa hai cấu hình. GPU lúc đo chạy ở xung trung vị 780 đến 802 MHz (tối đa 3\,105), 70 đến 72\,\textdegree C; số tuyệt đối ghi kèm điều kiện này, hai cấu hình đo cùng điều kiện.

\begin{figure}[!htbp]
\centering
% (a) A/B mù, lượt cả khung (bảng t8); (b) p95 ms nhịp thật có "Ảnh rõ" (bảng t9)
\begin{subfigure}[b]{0.47\textwidth}\centering
\begin{tikzpicture}
\begin{axis}[baocao, width=6.1cm, height=4.9cm, xbar, /pgf/bar width=12pt,
  ymajorgrids=false, xmajorgrids, xmin=0, xmax=1.18, xlabel={tỉ lệ lượt bên thứ nhất thắng},
  xtick={0,0.25,0.5,0.75,1},
  ymin=0.4, ymax=3.6, ytick={1,2,3},
  yticklabels={{Real-ESRGAN $\times$2\\vs cổ điển},{Real-ESRGAN $\times$2\\vs gốc},{cổ điển\\vs gốc}},
  y tick label style={align=right,font=\scriptsize}, /pgf/number format/precision=2,
  error bars/x dir=both, error bars/x explicit, error bars/error bar style={line width=0.7pt,black!70},
  error bars/error mark options={mark size=3pt,black!70}]
\addplot[fill=cXanh!85, draw=none, bar shift=0pt] coordinates {(1.0,3) -= (0.133,0)};
\addplot[fill=cTim!80, draw=none, bar shift=0pt] coordinates {(1.0,2) -= (0.133,0) (0.64,1) += (0.158,0) -= (0.195,0)};
\draw[cXam, dashed] (axis cs:0.5,0.4) -- (axis cs:0.5,3.6);
\node[font=\scriptsize, anchor=west, white] at (axis cs:0.02,3) {25/25};
\node[font=\scriptsize, anchor=west, white] at (axis cs:0.02,2) {25/25};
\node[font=\scriptsize, anchor=west, white] at (axis cs:0.02,1) {16/25};
\end{axis}
\end{tikzpicture}
\caption{Thử mù A/B, 25 cặp mỗi so sánh, CI Wilson 95\,\%}
\label{fig:ab}
\end{subfigure}\hfill
\begin{subfigure}[b]{0.5\textwidth}\centering
\begin{tikzpicture}
\begin{axis}[baocao, width=7.4cm, height=4.9cm, ybar=1pt, bar width=9pt,
  ymin=0, ymax=140, ylabel={p95 ms, nhịp thật},
  symbolic x coords={det,ver,enh,tong}, xtick={det,ver,enh,tong},
  xticklabels={detector,verifier,ảnh rõ,tổng},
  x tick label style={font=\scriptsize}, enlarge x limits=0.18,
  legend style={at={(0.03,0.97)},anchor=north west,legend columns=1},
  nodes near coords, every node near coord/.append style={font=\tiny, /pgf/number format/precision=1},
  /pgf/number format/precision=1]
\addplot[fill=cXam!60, draw=none] coordinates {(det,63.3) (enh,34.4) (tong,99.1)}; \addlegendentry{một tầng + NAFNet}
\addplot[fill=cXanh!85, draw=none] coordinates {(det,63.4) (ver,8.9) (enh,19.3) (tong,89.5)}; \addlegendentry{ba tầng + bộ lọc cổ điển}
\end{axis}
\end{tikzpicture}
\caption{IOT-CV phát lại 600 khung, có người xem ``Ảnh rõ''}
\label{fig:ms}
\end{subfigure}

\caption[Thử mù A/B và thời gian xử lý trong IOT-CV]{(a) Thử mù A/B, lượt cả khung; vạch đứt là 0,5. (b) p95 từng bước trong IOT-CV ở hai cấu hình: nhận diện một tầng + NAFNet và ba tầng + bộ lọc cổ điển.}
\label{fig:ab-ms}
\end{figure}
